\documentclass[11pt]{article}
\usepackage[margin=1in]{geometry}
\usepackage{amsmath,amssymb,amsthm}
\usepackage{booktabs}
\usepackage{array}
\usepackage{microtype}
\usepackage{xcolor}
\usepackage{hyperref}
\hypersetup{colorlinks=true,linkcolor=blue!50!black,citecolor=blue!50!black,urlcolor=blue!50!black}

\newtheorem{proposition}{Proposition}
\newtheorem{lemma}{Lemma}
\theoremstyle{remark}
\newtheorem{remark}{Remark}

\newcommand{\F}{\mathbb{F}}
\newcommand{\R}{\mathbb{R}}
\newcommand{\Aff}{\mathrm{Aff}}
\newcommand{\relu}{\mathrm{ReLU}}
\newcommand{\cost}{C}
\newcommand{\gam}{\gamma}
\newcommand{\abs}[1]{\lvert #1 \rvert}
\newcommand{\norm}[1]{\lVert #1 \rVert}
\newcommand{\ind}[1]{\mathbf{1}[#1]}

\title{Margins of coset and clock circuits on $\Aff(\F_p)$}
\author{Jaeho Lee}
\date{October 2026}

\begin{document}
\maketitle

\begin{abstract}
A one-hidden-layer ReLU network trained with weight decay on multiplication in the affine group $\Aff(\F_p)$ organises its
neurons by cosets of point stabilisers (the ``plain'' circuit) rather than by the Fourier basis of the induced representation
(the ``twisted'' or clock circuit). Weight decay selects for margin per unit weight norm, so the question is which circuit has
the larger normalised margin. We compute the margin and the norm cost of both hand-built circuits in closed form, verify every
formula by brute force at $p = 7, 11, 13$, and find: the plain circuit's normalised margin scales as $p^{-7/2}$, the clock's as
$p^{-3}$, but with a constant $23.6/\sqrt{p}$ smaller, so the clock mechanism would overtake the plain one only near $p \approx 555$,
and the complete clock circuit never does because the unit part it needs is itself a coordinate circuit. The gap factorises into
the ReLU attenuation of the clock's product term ($3\pi/2$), its worst-case input factor ($3\pi/4$), the cost of silencing a neuron
off its fibre ($\sqrt{3}$), and the plain circuit's own gap-to-norm handicap ($1.22/\sqrt{p}$), the only $p$-dependent factor.
\end{abstract}

\section{Setting}

\paragraph{Group.} $G = \Aff(\F_p) = \{x : j \mapsto a_x j + b_x \mid a_x \in \F_p^{*},\, b_x \in \F_p\}$, $\abs{G} = p(p-1)$.
Composition is $(xy)(j) = x(y(j))$, so $a_{xy} = a_x a_y$ and $b_{xy} = a_x b_y + b_x$. We write $x^{-1}(t) = (t - b_x)/a_x$ for the
point that $x$ sends to $t$.

\paragraph{Network.} Hidden neuron $k$ has pre-activation $L_k(x) + R_k(y)$ with $L_k, R_k : G \to \R$ arbitrary (one-hot embeddings
make every function available, and a constant plays the role of a bias). The logit of output $z$ is
\[
f_z(x,y) = \sum_k w_k(z)\, \relu\bigl(L_k(x) + R_k(y)\bigr), \qquad w_k : G \to \R .
\]

\paragraph{Margin and cost.} The margin is $\gam = \min_{(x,y)} \min_{z \ne xy} \bigl(f_{xy}(x,y) - f_z(x,y)\bigr)$. For a two-layer
ReLU network with $L^2$ weight decay, once the in and out norms of each neuron balance, the penalty equals
\[
\cost = \sum_k \norm{(L_k, R_k)}_2 \, \norm{w_k}_2, \qquad \norm{(L_k,R_k)}_2^2 = \sum_{g \in G} L_k(g)^2 + \sum_{g \in G} R_k(g)^2 .
\]
The ratio $\gam/\cost$ is scale invariant and is the quantity that margin maximisation under weight decay drives up
(Morwani et al.\ 2023). We compare it between circuits.

\paragraph{Notation.} For a profile $F : \F_p \to \R$ let $g_F$ be its smallest gap between distinct values and
$S_F = \sum_j F(j)^2$; likewise $g_f$, $S_f$ for a profile $f$ on $\F_p^{*}$. With the evenly spaced centred choices
$F(j) = j - \tfrac{p-1}{2}$ and $f$ taking the values $\{-\tfrac{p-2}{2}, \dots, \tfrac{p-2}{2}\}$ in some fixed order,
\[
g_F = g_f = 1, \qquad S_F = \frac{p(p^2-1)}{12}, \qquad S_f = \frac{p(p-1)(p-2)}{12},
\]
and we abbreviate $a_p = \sqrt{2 S_F} = \sqrt{p(p^2-1)/6}$, $b_p = \sqrt{2 S_f} = \sqrt{p(p-1)(p-2)/6}$.

\section{The plain coordinate circuit}

\subsection{Construction}

Fix a set $J \subseteq \F_p$ of \emph{coordinates}. For every $j_0 \in J$ and every value $c \in \F_p$ there is a pair of neurons
\[
L(x) = \pm F\bigl(x^{-1}(c)\bigr), \qquad R(y) = \mp F\bigl(y(j_0)\bigr), \qquad w(z) = -\ind{z(j_0) = c},
\]
so the pair contributes $-\abs{F(x^{-1}(c)) - F(y(j_0))}$ to the logits of every $z$ with $z(j_0) = c$. Since $x^{-1}(c) = y(j_0)$
exactly when $(xy)(j_0) = c$, the correct output receives zero from every pair. The \emph{unit part} fixes $a_z$: for every
$v \in \F_p^{*}$ a pair with $L(x) = \pm f(a_x)$, $R(y) = \mp f(v / a_y)$, $w(z) = -\ind{a_z = v}$. The circuit has $2p\abs{J} + 2(p-1)$
neurons.

\subsection{The logit gap}

For $z \ne xy$,
\begin{equation}\label{eq:gap}
\Delta(x,y,z) = f_{xy}(x,y) - f_z(x,y) = \sum_{j_0 \in J} \abs{F\bigl(x^{-1}(z(j_0))\bigr) - F\bigl(y(j_0)\bigr)}
 + \abs{f(a_x) - f(a_z / a_y)},
\end{equation}
and the key identity, from $x^{-1}(t) = (t - b_x)/a_x$, is
\begin{equation}\label{eq:shift}
x^{-1}\bigl(z(j_0)\bigr) - y(j_0) = \frac{z(j_0) - (xy)(j_0)}{a_x} = \frac{(a_z - a_{xy})\, j_0 + (b_z - b_{xy})}{a_x} \pmod p .
\end{equation}
With $F(j) = j - \tfrac{p-1}{2}$ each coordinate term is the ordinary integer distance between representatives in
$\{0, \dots, p-1\}$, not the cyclic distance; this is what produces the edge effect below.

\subsection{Exact margin}

\begin{proposition}\label{prop:plain}
With evenly spaced $F$ and $f$ at the same scale,
\begin{align*}
\gam(J) &= \min\bigl(\abs{J}\, g_F,\; (\abs{J}-1)\, g_F + g_f\bigr) = \abs{J} & (1 \le \abs{J} \le p-1), \\
\gam(J) &= 2(p-1)\, g_F = 2(p-1) & (\abs{J} = p,\ p \ge 7).
\end{align*}
\end{proposition}

\begin{proof}[Proof sketch]
\emph{Case A: $a_z = a_{xy}$, $b_z \ne b_{xy}$.} The unit term vanishes and by \eqref{eq:shift} the shift
$\delta = (b_z - b_{xy})/a_x$ is the same for every $j_0$. Writing $\delta' \in \{1, \dots, p-1\}$, the coordinate term is $\delta'$
when $y(j_0) + \delta' \le p-1$ and $p - \delta'$ when it wraps. The minimum is $\delta' = 1$ (or $p-1$) with $y(j_0) \ne p-1$ for all
$j_0 \in J$, which is possible iff $\abs{J} \le p-1$, giving $\abs{J} g_F$. For $\abs{J} = p$ exactly one coordinate wraps, and the
total is $(p-1) g_F + (p-1) g_F = 2(p-1) g_F$; any other $\delta'$ gives $\sum_r \abs{F(r+\delta') - F(r)} \ge 2(\max F - \min F) = 2(p-1) g_F$
over a full cycle, so $\delta' = \pm 1$ is the worst.

\emph{Case B: $a_z \ne a_{xy}$.} The unit term is at least $g_f$. The maps $z$ and $xy$ agree at exactly one point
$j^{*} = (b_{xy} - b_z)/(a_z - a_{xy})$, so every other coordinate contributes at least $g_F$ and
$\Delta \ge (\abs{J} - \ind{j^{*} \in J})\, g_F + g_f \ge (\abs{J}-1) g_F + g_f$, which is attained. For $\abs{J} = p$ the shift in
\eqref{eq:shift} runs over all of $\F_p^{*}$ as $j_0$ runs over $\F_p \setminus \{j^{*}\}$, so
$\Delta \ge g_F \sum_{d=1}^{p-1} \min(d, p-d) + g_f = g_F\,\tfrac{p^2-1}{4} + g_f$, which exceeds $2(p-1) g_F$ for $p \ge 7$
(brute force: $13$ vs $12$ at $p=7$, $31$ vs $20$ at $p=11$, $43$ vs $24$ at $p=13$).
\end{proof}

\begin{remark}
The all-coordinates margins $20$ ($p=11$) and $24$ ($p=13$) are $2(p-1)$, not $p-1$ or $p$: the worst case is case A with
$x = y = \mathrm{id}$ and $z$ a translation by one, which disagrees with $xy$ by one step on every coordinate, and the coordinate
where $y(j_0) = p-1$ pays $\abs{F(0) - F(p-1)} = p-1$ instead of $1$ because $F$ is a line, not a circle. The last coordinate
therefore doubles the margin; this is a property of the one-dimensional profile, not of the circuit.
\end{remark}

\begin{remark}
For $\abs{J} \ge 2$ the unit part is not needed for correctness (two coordinates determine an affine map), and
$\gam = (\abs{J}-1) g_F$ without it. At $\abs{J} = p$ it is dead weight: the margin is unchanged and the cost drops, so
$\gam/\cost$ rises from $0.517 \times 10^{-3}$ to $0.557 \times 10^{-3}$ at $p = 11$.
\end{remark}

\subsection{Exact cost}

A coordinate neuron's side functions $x \mapsto F(x^{-1}(c))$ and $y \mapsto F(y(j_0))$ take each value of $F$ on exactly $p-1$
group elements, so $\norm{L}^2 = \norm{R}^2 = (p-1) S_F$, and $w$ is the indicator of a set of size $p-1$. A unit neuron's sides take
each value of $f$ on $p$ elements and its readout indicates a set of size $p$. Hence
\begin{equation}\label{eq:cost}
\cost(J) = 2p\abs{J}\,(p-1)\,a_p\, g_F + 2(p-1)\, p\, b_p\, g_f = 2p(p-1)\bigl[\abs{J}\, a_p + b_p\bigr] \quad (g_F = g_f = 1).
\end{equation}

\subsection{Normalised margin and its scaling}

\begin{equation}\label{eq:plainratio}
\frac{\gam}{\cost} = \frac{\abs{J}}{2p(p-1)\,[\abs{J}\, a_p + b_p]} \quad (\abs{J} \le p-1), \qquad
\frac{\gam}{\cost} = \frac{1}{p\,[p\, a_p + b_p]} \quad (\abs{J} = p).
\end{equation}

\begin{table}[h]\centering
\begin{tabular}{c c c c c c}
\toprule
$p$ & $\abs{J}=1$ & $\abs{J}=2$ & $\abs{J}=\lfloor p/2 \rfloor$ & $\abs{J}=p-1$ & $\abs{J}=p$ \\
\midrule
7  & 0.889 & 1.140 & 1.259 & 1.406 & 2.450 \\
11 & 0.164 & 0.214 & 0.261 & 0.282 & 0.517 \\
13 & 0.089 & 0.116 & 0.146 & 0.156 & 0.290 \\
\bottomrule
\end{tabular}
\caption{$\gam/\cost$ in units of $10^{-3}$ for the plain circuit with evenly spaced $F$ and $f$. All entries agree with the
brute-force evaluation of the circuit on every input pair; \eqref{eq:cost} reproduces the brute-force cost to $10^{-10}$.}
\end{table}

Since $a_p, b_p \sim p^{3/2}/\sqrt{6}$, every $\abs{J}$ scales as $p^{-7/2}$ and only the constant moves:
$\tfrac{\sqrt{6}}{4}$ at $\abs{J} = 1$, $\tfrac{\sqrt{6}}{2}$ at $\abs{J} = p-1$, $\sqrt{6}$ at $\abs{J} = p$. The exponent is
(number of neuron pairs $p^2$) $\times$ (per-neuron norm $p^{5/2}$) $/$ (margin $2p$).

\subsection{The optimal profile}

\begin{lemma}
Among centred injective $F$ with smallest gap $g_F$, the evenly spaced profile minimises $S_F$, hence maximises $g_F / \norm{F}$.
\end{lemma}
\begin{proof}
For centred $F$ with sorted values $v_1 < \dots < v_p$, $S_F = \tfrac{1}{p}\sum_{i<k}(v_k - v_i)^2 \ge \tfrac{g_F^2}{p}\sum_{i<k}(k-i)^2
= g_F^2\, \tfrac{p(p^2-1)}{12}$, with equality iff the values are evenly spaced with gap $g_F$.
\end{proof}

For $\abs{J} = 1$ the margin is exactly $\min(g_F, g_f)$, so the ratio $\min(g_F,g_f)/(2p(p-1)[a_p g_F + b_p g_f])$ is maximised at
$g_F = g_f$: any imbalance adds cost without margin. The same holds for $2 \le \abs{J} \le p-1$, because the kink of
$\min(\abs{J} g_F, (\abs{J}-1) g_F + g_f)$ is the optimum when $b_p / a_p = \sqrt{(p-2)/(p+1)} < \abs{J}/(\abs{J}-1)$; for $\abs{J} = p$
the optimum is $g_f = 0$. The structural inefficiency of the plain circuit is now visible: $\norm{F} \sim p^{3/2} g_F$ while only
the smallest gap enters the margin.

\section{The twisted clock circuit}

\subsection{Construction and separability}

The $(p-1)$-dimensional irrep of $G$ is induced from a non-trivial character of the translation subgroup, so its ``twisted'' basis
functions are supported on one fibre $a_x = \alpha$ and are a single Fourier mode in $b_x$. On that fibre the translation part of
the product is $b_{xy} = b_x + \alpha b_y$: modular addition with $y$ rescaled by $\alpha$. The clock circuit has, for every fibre
$\alpha \in \F_p^{*}$, frequency $k$ in a set $\mathcal{F} \subseteq \{1, \dots, \tfrac{p-1}{2}\}$ and phase $\phi_i = 2\pi i / K$,
one neuron
\[
L(x) = \begin{cases} \cos(u + \phi_i) & a_x = \alpha \\ -1 & a_x \ne \alpha \end{cases}, \quad
R(y) = \cos(v + \phi_i), \quad w(z) = \cos(\theta + 2\phi_i),
\]
with $u = 2\pi k b_x / p$, $v = 2\pi k \alpha b_y / p$, $\theta = 2\pi k b_z / p$, so that $u + v = 2\pi k b_{xy}/p$. Off the fibre
$L + R \le 0$ and the ReLU is silent. The same unit part as before fixes $a_z$. Because the readout depends on $b_z$ only and the
unit part on $a_z$ only, $f_z = \Lambda(b_z) + U(a_z)$ and
\[
\gam = \min(\gam_{\mathrm{clock}}, \gam_{\mathrm{unit}}), \qquad \gam_{\mathrm{unit}} = g_f = 1,
\]
where $\gam_{\mathrm{clock}}$ is the margin over wrong $z$ with $a_z = a_{xy}$. Every fibre sees the same structure, so the clock
margin can be computed on one fibre.

\subsection{The phase sum}

Write $\relu(\cos(u+\phi) + \cos(v+\phi)) = 2\abs{c}\,\relu(\cos(\psi + \phi + \pi s))$ with $c = \cos\tfrac{u-v}{2}$,
$\psi = \tfrac{u+v}{2}$, $s = \ind{c < 0}$. The Fourier series of the rectified cosine is
\[
\relu(\cos t) = \sum_{n \ge 0} a_n \cos(nt), \qquad a_0 = \tfrac{1}{\pi},\quad a_1 = \tfrac12,\quad
a_{2m} = \frac{2}{\pi}\,\frac{(-1)^{m+1}}{4m^2 - 1},\quad a_{2m+1} = 0\ (m \ge 1),
\]
so $a_2 = \tfrac{2}{3\pi}$, $a_4 = -\tfrac{2}{15\pi}$, $a_6 = \tfrac{2}{35\pi}$. The phase sum identity
\[
\sum_{i=0}^{K-1} \cos(n\phi_i + A)\cos(2\phi_i + B) = \frac{K}{2}\Bigl[\cos(A-B)\,\ind{n \equiv 2 \bmod K} + \cos(A+B)\,\ind{n \equiv -2 \bmod K}\Bigr]
\]
gives the exact finite-$K$ result
\begin{equation}\label{eq:SK}
\begin{split}
S_K(u,v,\theta) &= \sum_i \relu\bigl(\cos(u+\phi_i)+\cos(v+\phi_i)\bigr)\cos(\theta + 2\phi_i) \\
&= K \abs{\cos\tfrac{u-v}{2}} \sum_{m \ge 1} a_{2m}\Bigl[\cos\bigl(m(u+v) - \theta\bigr)\,\ind{2m \equiv 2 \bmod K} \\
&\hspace{9em} + \cos\bigl(m(u+v) + \theta\bigr)\,\ind{2m \equiv -2 \bmod K}\Bigr]
\end{split}
\end{equation}
(plus an $a_1$ term if $K \mid 3$ and an $a_0$ term if $K \mid 2$), the sign $s$ dropping out of the $n = 2$ term. For $K \to \infty$ only
$m = 1$ survives:
\begin{equation}\label{eq:Sinf}
S_\infty(u,v,\theta) = c_K \abs{\cos\tfrac{u-v}{2}}\cos(u + v - \theta), \qquad c_K = K a_2 = \frac{2K}{3\pi} = 0.2122\,K .
\end{equation}
The main term has exactly this coefficient for every $K \ge 5$; the only finite-$K$ effects are the aliased harmonics. The factor
$\abs{\cos\tfrac{u-v}{2}} = \abs{\cos(\pi k(b_x - \alpha b_y)/p)}$ depends on the input and on $k$, so margins are taken at the worst
input.

\subsection{Aliasing}

Harmonic $n$ of $\relu(\cos t)$ reaches the readout iff $n \equiv \pm 2 \pmod K$.

\begin{table}[h]\centering
\begin{tabular}{c l l l}
\toprule
$K$ & surviving harmonics & lowest alias, size relative to $a_2$ & outcome \\
\midrule
3  & $2$; $1, 4, 8, 10, \dots$ & $n=1$: $a_1/a_2 = 3\pi/4 = 2.36$ & fails \\
4  & $2$ in both channels; $6, 10, \dots$ & the main term aliases with itself & fails \\
6  & $2$; $4, 8, 10, \dots$ & $n=4$: $1/5$ & works, marginal \\
8  & $2$; $6, 10, 14, \dots$ & $n=6$: $3/35 = 0.086$ & works \\
12 & $2$; $10, 14, 22, \dots$ & $n=10$: $3/99 = 0.030$ & works \\
\bottomrule
\end{tabular}
\end{table}

At $K = 4$ both channels of $n = 2$ survive and $S_4 = 8\abs{c}\sum_{m\ \mathrm{odd}} a_{2m}\cos(m(u+v))\cos\theta$: the logit is a
function of $u+v$ times $\cos\theta$, the argmax over $b_z$ can only be $0$ or $(p \pm 1)/2$, and the accuracy is $0.18$ at $p = 11$.
For even $K \ge 6$ the lowest alias is $n = K - 2$ with relative size $3/((K-2)^2 - 1)$. An alias of relative size $\varepsilon$ damages
a single-frequency margin by about $\varepsilon p / \pi$ (its variation between adjacent $b_z$ is $2\varepsilon\sin(\pi/p)$ against a main-term
gap of $1 - \cos(2\pi/p)$): $55\%$, $21\%$, $6\%$ at $K = 6, 8, 12$ for $p = 11$. With all frequencies the gap is $O(p)$ and the aliases
cost only $8\%$ ($K=8$) and $2.5\%$ ($K=12$).

\subsection{One frequency}

On a fibre with $m = b_x - \alpha b_y$, $\Lambda(b_z) = c_K \abs{\cos(\pi k m / p)}\cos(2\pi k(b_{xy} - b_z)/p)$ (alias-free). The gap to the
nearest wrong $b_z$ is $c_K\abs{\cos(\pi k m/p)}(1 - \cos(2\pi/p))$, and the worst input has $km \equiv (p \pm 1)/2$, where the factor is
$\cos\tfrac{\pi(p-1)}{2p} = \sin\tfrac{\pi}{2p}$. Hence
\begin{equation}
\gam_{\mathrm{clock}}^{(1)} = c_K\,\bigl(1 - \cos\tfrac{2\pi}{p}\bigr)\sin\tfrac{\pi}{2p} \approx \frac{2K\pi^2}{3p^3}.
\end{equation}
Per neuron, $\norm{L}^2 = p(p-2) + p/2$ (the silencing constant sits on the $p(p-2)$ off-fibre elements), $\norm{R}^2 = \norm{w}^2 = p(p-1)/2$,
so the product of norms is $\tfrac{p}{2}\sqrt{(3p-4)(p-1)}$, and with $(p-1)K\abs{\mathcal{F}}$ neurons
\begin{equation}
\cost_{\mathrm{clock}} = (p-1)\,\abs{\mathcal{F}}\,K\,\frac{p}{2}\sqrt{(3p-4)(p-1)},
\end{equation}
and for one frequency, with the unit cost dropped,
\begin{equation}
\frac{\gam}{\cost} \approx \frac{4\pi^2}{3\sqrt{3}}\,p^{-6} = 7.6\,p^{-6}.
\end{equation}

\subsection{All frequencies}

With $\mathcal{F} = \{1, \dots, \tfrac{p-1}{2}\}$, $\Lambda(b_z) = c_K\sum_k \abs{\cos(\pi k m/p)}\cos(2\pi k d/p)$ with $d = b_{xy} - b_z$.
If $m = 0$ all factors are $1$ and the Dirichlet kernel $\sum_k \cos(2\pi k d/p) = (p\,\delta_{d,0} - 1)/2$ gives a gap of $c_K p/2$. If
$m \ne 0$ the factors $\{\abs{\cos(\pi k m/p)}\}_k$ are a permutation of $\{\cos(\pi j/p)\}_{j=1}^{(p-1)/2}$, so with $e = d/m$ the gap is
$c_K\, T(e)$, $T(e) = \sum_{j=1}^{(p-1)/2}\cos(\pi j/p)\bigl(1 - \cos(2\pi j e/p)\bigr)$, minimised at $e = \pm 1$ for every $p$ tested, with
\begin{equation}
T(1) = \frac14\Bigl[\csc\frac{\pi}{2p} + \csc\frac{3\pi}{2p}\Bigr] = \frac{2p}{3\pi} + O(p^{-1}), \qquad
\gam_{\mathrm{clock}}^{(\mathrm{all})} = c_K\, T(1) \approx \frac{4Kp}{9\pi^2} = 0.045\,Kp .
\end{equation}
The effective worst-case attenuation is therefore $T(1)/(p/2) \to 4/(3\pi) = 0.42$, a constant, rather than the $\sin(\pi/2p)$ of a single
frequency. The margin of the circuit as built is $\min(\gam_{\mathrm{clock}}, 1)$, and since $\gam_{\mathrm{clock}} > 1$ at $p \ge 7$, the
reported margin $1.000$ is the unit part's. Rebalancing the two parts (scaling each to margin $1$) gives
\begin{equation}
\Bigl(\frac{\gam}{\cost}\Bigr)_{\mathrm{rebalanced}} = \frac{1}{\cost_{\mathrm{clock}}/\gam_{\mathrm{clock}} + \cost_{\mathrm{unit}}/\gam_{\mathrm{unit}}}.
\end{equation}

\begin{table}[h]\centering
\begin{tabular}{c c c c c c c c c}
\toprule
$p$ & $K$ & freqs & $\gam$ & $\gam_{\mathrm{clock}}$ (alias-free) & $\cost_{\mathrm{clock}}$ & $\cost_{\mathrm{unit}}$ & $\gam/\cost$ as built & rebalanced \\
\midrule
7  & 4  & 1   & $-0.35$ & $-0.35$ (0.071) & 848 & 497 & fails & fails \\
7  & 6  & 1   & 0.067 & 0.067 (0.107) & 1273 & 497 & 0.038 & 0.051 \\
7  & 8  & 1   & 0.122 & 0.122 (0.142) & 1697 & 497 & 0.055 & 0.069 \\
7  & 8  & 3   & 1.000 & 2.387 (2.588) & 5090 & 497 & 0.179 & 0.380 \\
11 & 4  & 1   & $-0.35$ & $-0.35$ (0.019) & 3747 & 2826 & fails & fails \\
11 & 6  & 1   & 0.013 & 0.013 (0.029) & 5620 & 2826 & 0.0015 & 0.0023 \\
11 & 8  & 1   & 0.030 & 0.030 (0.038) & 7493 & 2826 & 0.0029 & 0.0040 \\
11 & 12 & 1   & 0.054 & 0.054 (0.058) & 11239 & 2826 & 0.0038 & 0.0047 \\
11 & 8  & 5   & 1.000 & 3.694 (4.004) & 37465 & 2826 & 0.025 & 0.077 \\
11 & 12 & 5   & 1.000 & 5.855 (6.006) & 56197 & 2826 & 0.017 & 0.081 \\
13 & 8  & 1   & 0.018 & 0.018 (0.023) & 12788 & 5276 & 0.0010 & 0.0014 \\
13 & 8  & 6   & 1.000 & 4.353 (4.718) & 76729 & 5276 & 0.012 & 0.044 \\
13 & 12 & 6   & 1.000 & 6.899 (7.077) & 115094 & 5276 & 0.008 & 0.046 \\
\bottomrule
\end{tabular}
\caption{Twisted clock circuit, brute force over all pairs; $\gam/\cost$ in units of $10^{-3}$. The closed forms reproduce every
entry to the digits shown.}
\end{table}

\subsection{Scaling and where the gap comes from}

Cost per unit of margin at leading order ($K$ cancels in the clock; alias-free):
\begin{align*}
\text{plain coordinates, } \abs{J} = p: &\quad \cost/\gam = p^2 a_p \sim \tfrac{1}{\sqrt6}\, p^{7/2} = 0.408\, p^{3.5}, \\
\text{clock part, all frequencies:} &\quad \cost_{\mathrm{clock}}/\gam_{\mathrm{clock}} \sim \tfrac{9\sqrt3\,\pi^2}{16}\, p^3 = 9.62\, p^{3}, \\
\text{unit part (both circuits):} &\quad \cost_{\mathrm{unit}}/\gam_{\mathrm{unit}} = 2p(p-1)b_p \sim \tfrac{2}{\sqrt6}\, p^{7/2} = 0.816\, p^{3.5}, \\
\text{clock part, one frequency:} &\quad \cost_{\mathrm{clock}}/\gam_{\mathrm{clock}} \sim \tfrac{3\sqrt3}{4\pi^2}\, p^6 = 0.13\, p^{6}.
\end{align*}
Hence $\gam/\cost \sim \sqrt6\, p^{-7/2}$ for the plain circuit, $\sim 0.104\, p^{-3}$ for the clock mechanism alone, and for the
complete rebalanced twisted circuit $1/(9.62\,p^3 + 0.816\,p^{3.5}) \to \tfrac{\sqrt6}{2}\, p^{-7/2}$, half of the plain value, because
the unit part it needs is itself a coordinate circuit on $\F_p^{*}$. The clock mechanism overtakes the plain coordinates only at
$p \approx (9.62/0.408)^2 = 555$.

\begin{table}[h]\centering
\begin{tabular}{r c c c c}
\toprule
$p$ & plain $\abs{J}=p$ & twisted as built ($K=8$) & twisted rebalanced & clock part alone \\
\midrule
11   & $5.2\times10^{-4}$  & $2.5\times10^{-5}$  & $8.2\times10^{-5}$  & $1.1\times10^{-4}$ \\
13   & $2.9\times10^{-4}$  & $1.2\times10^{-5}$  & $4.6\times10^{-5}$  & $6.2\times10^{-5}$ \\
31   & $1.4\times10^{-5}$  & $3.3\times10^{-7}$  & $2.6\times10^{-6}$  & $3.9\times10^{-6}$ \\
101  & $2.3\times10^{-7}$  & $2.8\times10^{-9}$  & $5.6\times10^{-8}$  & $1.0\times10^{-7}$ \\
547  & $6.4\times10^{-10}$ & $3.2\times10^{-12}$ & $2.1\times10^{-10}$ & $6.4\times10^{-10}$ \\
\bottomrule
\end{tabular}
\caption{Leading-order $\gam/\cost$ of the four mechanisms. The plain circuit leads by $21\times$ and $24\times$ as built at
$p = 11, 13$, by $6.3\times$ after rebalancing.}
\end{table}

The ratio of the clock's cost per margin to the plain coordinates' factorises as
\[
\underbrace{\tfrac{3\pi}{2}}_{\text{ReLU attenuation}}\ \times\ \underbrace{\tfrac{3\pi}{4}}_{\text{worst-case } \abs{\cos}}\ \times\
\underbrace{\sqrt3}_{\text{silencing}}\ \times\ \underbrace{\tfrac{\sqrt6}{2}\, p^{-1/2}}_{\text{bounded features vs } \norm{F}/g_F \sim p^{3/2}}
\ =\ \frac{23.6}{\sqrt p},
\]
with finite-$K$ aliasing adding $1.08$ at $K = 8$. The ReLU attenuation is the fact that only the second $\phi$-harmonic, of weight
$a_2 = 2/(3\pi)$, reaches the readout; the worst-case factor is $T(1)/(p/2) \to 4/(3\pi)$; silencing raises $\norm{(L,R)}^2$ from
$p^2/2$ to $p(3p-4)/2$; and the last factor is the plain circuit's handicap, bounded clock features against a profile whose norm grows
like $p^{3/2}$ while only its smallest gap counts. The per-fibre duplication ($p-1$ copies) is matched one for one by the plain
circuit's $p$ value copies per coordinate, so both circuits have $\Theta(p^2)$ neurons and the duplication is not a relative factor;
what it does is force the silencing. At $p = 11$, $K = 8$ the as-built gap of $20.8$ decomposes as $3.1$ (unit-limited margin)
$\times\, 5.65$ (clock vs coordinate cost per margin: count $1.65$, per-neuron $0.63$, margin $5.0$) $\times$ the amortisation of the
unit part over a margin of $20$ rather than $3.7$.

\section{Cheaper silencing?}

\paragraph{A smaller offset fails.} With $L = -\eta$ off the fibre and $R$ unchanged, an inactive neuron contributes
$\sum_i \relu(\cos(v+\phi_i) - \eta)\cos(\theta + 2\phi_i) \approx \tfrac{K}{2}\, a_2(\eta)\cos(2v - \theta)$ with the exact coefficient
\[
a_2(\eta) = \frac{2}{3\pi}\,(1 - \eta^2)^{3/2},
\]
a spurious clock pointing at $b_z = 2\alpha b_y$. For $b_y = 0$ all $p-2$ inactive fibres are coherent and produce a Dirichlet spike of
height $\sim \tfrac{K}{4} a_2(\eta)\, p(p-2)$ on $b_z = 0$; keeping it below $\gam_{\mathrm{clock}} \sim 4Kp/(9\pi^2)$ requires
$1 - \eta^2 \lesssim (0.85/(p-2))^{2/3}$, so the admissible saving is $O(p^{-2/3})$ of the silencing cost. Brute force at $p = 11$:
$\eta = 0.95$ saves $3\%$ of $\cost_{\mathrm{clock}}$ and cuts the clock margin from $3.69$ to $1.50$; $\eta = 0.9$ gives a negative margin.

\paragraph{Splitting the constant works.} Put $L = \cos(u+\phi) + \beta$ on the fibre, $L = -(1-\beta)$ off it, $R = \cos(v+\phi) - \beta$.
On the fibre $L + R$ is unchanged, off it $L + R = \cos(v+\phi) - 1 \le 0$, so every activation and the margin are identical, while
\[
\norm{(L,R)}^2 = p(p-2)(1-\beta)^2 + p\bigl(\tfrac12 + \beta^2\bigr) + p(p-1)\bigl(\tfrac12 + \beta^2\bigr)
\]
is minimised at $\beta^{*} = \tfrac{p-2}{2(p-1)} \approx \tfrac12$ with value $\tfrac{p^2(2p-3)}{2(p-1)} \sim p^2$ against $\tfrac{p(3p-4)}{2} \sim \tfrac{3p^2}{2}$,
a cost factor $\to \sqrt{2/3} = 0.816$ ($0.849$ at $p = 11$, brute force $0.8489$). This is optimal within exact silencing of an exact clock:
the constant $1$ must be paid on about $p^2$ entries of $L$ or of $R$, and $(1-\beta)^2 + \beta^2 \ge \tfrac12$. It moves the clock's constant
to $0.127\, p^{-3}$ and the crossover to $p \approx 370$, not the exponents.

\section{The duplication cannot be avoided}

The twisted neuron needs its $y$-side feature at frequency $k a_x$, and $R(y)$ cannot see $a_x$; the only way to get $a_x$-dependence into
$L(x) + R(y)$ is one copy of $R$ per value of $a_x$, and each copy must be silenced off its fibre because the ReLU is nonlinear (a zero
$L$ off the fibre would still pass $\relu(R(y))$ through). Moving the twist to the $x$ side, $L = \cos(2\pi k b_x/(a_x p) + \phi)$,
$R = \cos(2\pi k b_y/p + \phi)$, gives $u + v = 2\pi k b_{xy}/(a_x p)$, which the readout would have to undo with $a_x$; but $w$ depends on
$z = xy$ only and $a_x$ is not a function of $z$. Representation-theoretically, the $(p-1)$-dimensional irrep is induced from a character
of the translation subgroup, so each of its matrix coefficients is supported on a single fibre: any neuron reading an irrep feature of
$x$ is fibre-local by construction, and the $p-1$ fibres are the $p-1$ rows of the representation matrix. The duplication is the price of
the irrep being induced, not of this particular circuit.

\section{What was verified}

Everything above was checked numerically. The checking scripts are \texttt{affine\_margins\_check.py} and
\texttt{affine\_margins\_split.py} in the repository, with their outputs under \texttt{results/affine}. Checked:
the Fourier coefficients of $\relu(\cos t)$; the exact alias series \eqref{eq:SK} against the direct phase sum for
$K \in \{3, 4, 6, 8, 12, 2000\}$; $a_2(\eta)$ and $T(1)$ in closed form and the minimiser $e = \pm 1$ for $p$ from $7$ to $101$;
both circuits built exactly as specified and
evaluated over all $\abs{G}^2$ input pairs and all outputs at $p = 7, 11, 13$, including the case A / case B minima separately, the
$\abs{J} = p$ minimiser, the no-unit-part and unequal-scale variants, the worst single-frequency input, the $\eta < 1$ and split-constant
variants, and the $K = 4$ collapse. Approximations appear only where labelled: alias-free ($K \to \infty$) clock margins, the case B lower
bound at $\abs{J} = p$ (valid for $p \ge 7$, checked at $7, 11, 13$), and the large-$p$ leading orders.

\end{document}
